% Ports: the rules \energeseport is meant to hold to.
%
% The golden images show where the pathways ended up; this asserts why. Each
% check raises a TeX error when it fails, so the suite reports a failure rather
% than a diagram that still looks plausible.
\documentclass{standalone}
\usepackage{energese}
\makeatletter

\newdimen\ts@x \newdimen\ts@y \newdimen\ts@u \newdimen\ts@v \newdimen\ts@tol
% Two points agree if they are within a printer's tolerance of each other: the
% port machinery runs through pgfmath, which rounds.
\ts@tol=0.2pt

\def\ts@anchor#1#2{\pgfpointanchor{#1}{#2}\ts@x=\pgf@x \ts@y=\pgf@y}
% The port just allocated, as a point. \energeseportname is `node.anchor'.
\def\ts@port{\expandafter\ts@portx\energeseportname\@nil}
\def\ts@portx#1.#2\@nil{\pgfpointanchor{#1}{#2}\ts@u=\pgf@x \ts@v=\pgf@y}

% Manhattan distance between the anchor point and the port point, left in
% \ts@u. Both assertions consume it.
\def\ts@dist{%
    \advance\ts@u by-\ts@x \advance\ts@v by-\ts@y
    \ifdim\ts@u<\z@ \ts@u=-\ts@u \fi
    \ifdim\ts@v<\z@ \ts@v=-\ts@v \fi
    \advance\ts@u by\ts@v}
\def\ts@same#1{\ts@dist
    \ifdim\ts@u>\ts@tol \PackageError{energese-test}{FAIL: #1}{}\fi}
\def\ts@apart#1{\ts@dist
    \ifdim\ts@u<\ts@tol \PackageError{energese-test}{FAIL: #1}{}\fi}

% Is the allocated port one of the ring's? Splitting the name on `ring' leaves
% nothing in front of it when it is.
\def\ts@onring#1{\expandafter\ts@onringx\energeseportname\@nil{#1}}
\def\ts@onringx#1.#2\@nil#3{\ts@onringy#2ring\@nil{#3}}
\def\ts@onringy#1ring#2\@nil#3{%
    \def\ts@head{#1}%
    \ifx\ts@head\@empty\else
        \PackageError{energese-test}{FAIL: #3}{}%
    \fi}

\begin{document}
\begin{tikzpicture}
    \node[energese consumer] (hub)   at ( 0, 0)   {Hub};
    \node[energese source]   (above) at (-3, 1.5) {A};
    \node[energese source]   (below) at (-3,-1.5) {B};
    \node[energese storage]  (right) at ( 3, 0)   {C};
    \node[energese text]     (note)  at ( 0,-3)   {note};

    \energeseresetports

    % 1. A named anchor on an axis is already a port, so quantising energy_in
    %    moves it nowhere: the ring always carries a port at 180 degrees.
    \energeseport{hub}{energy_in}{above}
    \ts@anchor{hub}{energy_in}\ts@port
    \ts@same{energy\string_in did not land on its own anchor}

    % 2. A port holds one pathway: the second request for the same anchor is
    %    given a different port ...
    \energeseport{hub}{energy_in}{below}
    \ts@anchor{hub}{energy_in}\ts@port
    \ts@apart{two pathways were given the same port}

    % 3. ... and it is one of this symbol's own ports, not a point invented
    %    for the occasion.
    \ts@onring{the second pathway was not put on the ring}

    % 4. A third is pushed further out again, rather than back onto the first.
    \energeseport{hub}{energy_in}{right}
    \ts@anchor{hub}{energy_in}\ts@port
    \ts@apart{the third pathway was given the anchor's own port}

    % 5. A shape with no ring keeps the anchor it was asked for and claims
    %    nothing: a `text' node takes its pathways at the centre.
    \energeseport{note}{center}{hub}
    \def\ts@expect{note.center}
    \ifx\energeseportname\ts@expect\else
        \PackageError{energese-test}{FAIL: a ringless shape was given a port}{}%
    \fi

    % 6. Claims last one diagram. After a reset the anchor's own port is free
    %    again and the next pathway gets it back.
    \energeseresetports
    \energeseport{hub}{energy_in}{above}
    \ts@anchor{hub}{energy_in}\ts@port
    \ts@same{a reset did not release the ports}
\end{tikzpicture}
\end{document}
